\section{Pre-Token Hook Configuration}\label{sec:hook-configuration}

\textit{Normative language per \href{https://tools.ietf.org/html/rfc2119}{RFC~2119}.}

\subsection{Purpose and Scope}

The pre-token hook is the mechanism by which the tenant's IdP injects the entitlement
claim into the access token before it is signed. It is the single most important
implementation step in the tenant's integration. Every behavioral requirement in
\hyperref[sec:entitlement-mapping]{\S4 Entitlement Mapping} depends on the hook being correctly configured.


This section defines:

\begin{enumerate}
  \item The four rules that ALL hooks MUST follow, regardless of IdP platform.
  \item Per-IdP implementation guidance for the four most common platforms.
  \item Auth0-specific operational notes addressing known failure modes.
\end{enumerate}

Tenants using an IdP not listed in \hyperref[sec:per-idp-implementation]{\S6.2} 
MUST implement a functionally equivalent mechanism and validate that the resulting 
token satisfies all requirements in 
\hyperref[sec:token-requirements]{\S5 Token Requirements} before onboarding can proceed.

\subsection{Per-IdP Implementation Table} \label{sec:per-idp-implementation}


\begin{table}[htbp]
\centering
\begin{tabularx}{\textwidth}{@{}lXp{4cm}@{}} % Slightly expanded note column width
\toprule
\textbf{IdP} & \textbf{Mechanism} & \textbf{Claim name note} \\
\midrule
AWS Cognito &
  Pre Token Generation Lambda trigger, event version \textbf{V2\_0} --- use
  \wrapcode{claimsAndScopeOverrideDetails.accessTokenGeneration.claimsToAddOrOverride}
  to set the claim. &
  Bare names (\texttt{entitlements}, \texttt{resources}) generally work.
  \texttt{cf:entitlements} is also supported. See Cognito \texttt{typ}/\texttt{token\_use}
  note in \hyperref[sec:token-type]{\S5.3.1}. \\[8pt]
Okta &
  Token Inline Hook --- return 
  \wrapcode{commands[].value.op = "add"} with
  \wrapcode{path = "/claims/cf:entitlements"}. &
  Bare names and namespaced names both work. \\[8pt]
Azure AD / Entra &
  Custom Claims Provider (custom authentication extension) configured as an
  \texttt{onTokenIssuanceStart} event handler. &
  Bare names work. Use attribute name matching the claim name in the extension's
  \wrapcode{additionalData} response. \\[8pt]
Auth0 &
  Post-Login \textbf{Action} ---
  \wrapcode{api.accessToken.setCustomClaim('cf:entitlements', [...])} &
  \textbf{MUST use \texttt{cf:entitlements}} --- see \hyperref[sec:auth0-notes]{\S6.3} for mandatory Auth0 notes. \\
\bottomrule
\end{tabularx}
\caption{Per-IdP hook implementation}
\end{table}

\subsection{Auth0-Specific Operational Notes}\label{sec:auth0-notes}

Auth0 introduces several behaviors that differ from other IdP platforms. Failure to
understand these is the primary source of Auth0 integration failures. The following
notes are drawn from direct operational experience and are mandatory reading for all
Auth0 tenants.

\subsubsection{Non-Namespaced Custom Claims Are Silently Dropped on Access Tokens}\label{sec:non-namespace-custom-claim}

Auth0 enforces a blocklist of non-namespaced custom claim names for access tokens. The
bare name \texttt{entitlements} is on this blocklist and will be silently dropped from
the access token, even if the Action correctly attempts to set it.

Auth0 tenants MUST use \texttt{cf:entitlements} as the claim name:

\begin{lstlisting}[language=js]
// Correct -- cf:entitlements is a namespaced claim and is not blocklisted
api.accessToken.setCustomClaim('cf:entitlements', canvasflowEntitlements);

// Incorrect -- 'entitlements' will be silently dropped by Auth0
api.accessToken.setCustomClaim('entitlements', canvasflowEntitlements);
\end{lstlisting}

Do not use \texttt{resources} or \texttt{entitlements} with Auth0. Neither is guaranteed
to survive the Auth0 access token pipeline.

\subsubsection{Action Must Be Deployed AND Attached in Triggers}\label{sec:action-deployment}

Auth0 separates the concepts of \emph{deploying} an Action and \emph{attaching} it to a
Trigger flow. An Action that has been written, saved, and deployed will do nothing until
it is also attached to the appropriate Trigger.

The correct procedure:
\begin{enumerate}
  \item Create and deploy the Post-Login Action in \textbf{Actions $\to$ Library}.
  \item Navigate to \textbf{Actions $\to$ Triggers} (the Auth0 dashboard may also show
    this as ``Flows'' in some UI versions).
  \item Select the \textbf{post-login} flow.
  \item Drag the deployed Action into the flow graph.
  \item Click \textbf{Apply}.
\end{enumerate}

An Action that is deployed but not attached produces no token output. There is no error
--- the token is simply issued without the \texttt{cf:entitlements} claim.

\subsubsection{\texttt{console.log} Does Not Appear in Monitoring $\to$ Logs}

\texttt{console.log} statements inside Auth0 Actions do NOT appear in
\textbf{Monitoring $\to$ Logs} in the Auth0 dashboard. Do not use \texttt{console.log}
to confirm that the hook executed or to inspect the entitlement values.

\textbf{How to verify the hook ran:} inspect the issued access token. Decode it at
\url{https://jwt.io} or use the Canvasflow OAuth Validator
(\url{https://oauth.canvasflow.work}). If \texttt{cf:entitlements} is present and
contains the expected values, the hook ran correctly. If the claim is absent, confirm
\hyperref[sec:non-namespace-custom-claim]{\S6.3.1} steps and 
\hyperref[sec:action-deployment]{\S6.3.2} above.


\subsubsection{API Registration Is Required for JWT Access Tokens}\label{sec:api-registration-required}

Auth0 issues opaque access tokens by default when only an Application is registered. To
obtain JWT access tokens, the tenant MUST register a separate \textbf{API} in the Auth0
dashboard. The API's \texttt{identifier} field becomes the \texttt{audience} parameter
required in every authorization request; it MUST be set to \texttt{canvasflow-api}.

The Application must then have the API authorized. The authorization request from the
PWA MUST include \texttt{audience=canvasflow-api}. Without both of these --- the API
registration and the audience parameter in the request --- Auth0 will issue an opaque
token that Canvasflow cannot validate.

Minimum scope for the token: \texttt{openid}. Additional scopes (\texttt{email},
\texttt{profile}) may be requested for ID token content but are not required for the
access token.

\pagebreak

\subsection{Hook Rules --- All IdPs}\label{sec:hook-rules}

Regardless of IdP platform, the pre-token hook MUST comply with all four of the
following rules:

\subsubsection{Rule~1: Re-evaluate on Every Invocation}

The hook MUST read the authenticated user's entitlements from the tenant's live source
of truth on every invocation --- both at initial login and at every token refresh. The
hook MUST NOT copy the entitlement claim from a prior token or from a prior hook result.

\textbf{Why.} The hook re-evaluation on refresh is the mechanism that keeps the effective
entitlement set current within the revocation window ($\leq$ 3600~seconds). If the hook
copies the previous token's claim rather than re-reading, a user whose access has been
revoked in the tenant's system will continue to receive premium content until the refresh
token is also revoked. This would undermine the security property stated in
\hyperref[sec:session-lifecycle]{Responsabilities \S3.2.6}.



\subsubsection{Rule~2: Map to Canvasflow Format Before Emitting}

The hook MUST map each internal tenant entitlement identifier to a Canvasflow resource
identifier in \texttt{cf:<type>:<id>} format before embedding it in the access token.
The hook SHOULD emit only values in this format.

Values without the \texttt{cf:} prefix are silently ignored by the Resource Server and
grant no access (see \hyperref[sec:mandatory-value-format]{Entitlement Mapping \S4.2}). A hook that emits only
non-prefixed values will produce a token that authenticates the user but grants no
premium content --- the same behavior as an empty entitlement array.

\subsubsection{Rule~3: Write into the Access Token, Not the ID Token}

The entitlement claim MUST be injected into the \textbf{access token}. It MUST NOT be
placed in the ID token.

This is the single most common misconfiguration in tenant integrations. The ID token is
consumed by the SPA for user identity display; it is never sent to the Canvasflow GraphQL
API. If the hook targets the ID token, the access token arrives at the Resource Server
with no entitlement claim, which is treated as free content only (see
\hyperref[sec:free-content-baseline]{Entitlement Mapping \S4.4}). The Canvasflow OAuth Validator (see
\hyperref[sec:tenant-side-setup]{On-Boarding \S9.3}) explicitly detects and flags this misconfiguration
before go-live.

\subsubsection{Rule~4: Emit an Explicit Empty Array for Free-Tier Users}

If the authenticated user has no premium entitlements, the hook MUST emit an explicit
empty array:

\begin{lstlisting}
"cf:entitlements": []
\end{lstlisting}

Do not omit the claim entirely when a user has no entitlements. An absent claim and an
empty array produce identical runtime behavior --- the user sees free content only --- but
an absent claim is logged by the Resource Server as a possible hook misconfiguration.
An explicit empty array signals to Canvasflow (and to the validator tool) that the hook
ran correctly and intentionally found no premium entitlements for this user.

\subsection{Hook Implementation Reference}

The following pseudocode illustrates the required hook logic for any IdP platform. The
implementation will differ in syntax and API; the logical structure MUST be preserved.

\begin{lstlisting}[language=js]
/**
 * Post-Login / Pre-Token Generation Hook
 * Implements token decoration rules for Canvasflow claims.
 */
function onPreTokenSign(user, event, context) {
    // Rule 1: Always read from live source of truth
    const tenantEntitlements = readUserEntitlements(user.id);

    // Rule 2: Map to Canvasflow cf:<type>:<id> format
    const canvasflowIds = [];
    
    tenantEntitlements.forEach(internalId => {
        const mapped = MAPPING_TABLE[internalId]; // tenant's own table
        
        if (mapped !== null && mapped !== undefined) {
            canvasflowIds.push(mapped); // e.g. "cf:issue:123"
        }
        // else: no mapping exists -- omit this value (it would be ignored anyway)
    });

    // Rule 4: Always emit the claim, even if empty
    // Rule 3: Target the access token, not the ID token
    // Note: Assuming `accessToken` is available in scope or extracted from context
    const accessToken = context.accessToken; 
    accessToken.setCustomClaim("cf:entitlements", canvasflowIds);

    return accessToken;
}
\end{lstlisting}

The \texttt{MAPPING\_TABLE} is the tenant's internal data structure mapping their product
or subscription identifiers to Canvasflow resource identifiers. It is maintained entirely
by the tenant.
